\section{Introduction}
\label{sec:intro}

Markov-chain ensemble methods have become the dominant tool for detecting
gerrymandering in U.S.\ congressional redistricting.  The core idea,
introduced by \citet{chikina2017assessing} and refined by \citet{deford2021recombination},
is to generate thousands of ``neutral'' maps by random walk on the space of
valid districting plans, then ask: where does the contested map fall in the
resulting distribution?  If the map lies in the 99th percentile of partisan
skew, one infers that geography alone cannot explain the skew, and that
intentional partisan sorting is the likely cause.

This inference engine is powerful precisely because it is comparative.
The ensemble does not tell you what the \emph{correct} map is.  It tells you
whether a given map is an outlier relative to what geography permits.  Courts
have begun to rely on this logic: ensemble evidence appeared in \case{League
of Women Voters of Pennsylvania v.\ Commonwealth of Pennsylvania} (2018),
\case{Common Cause v.\ Rucho} (2019), and was implicitly acknowledged in
\rucho{} even as the majority declined to apply a federal standard.

The \ar{} algorithm occupies a different position in this landscape.
Rather than sampling the space to characterise it, \ar{} directly computes
the minimum-edge-cut partition of the state's census-tract graph under the
prime factorization of the seat count.  This is not a sampled plan---it is
a \emph{local optimum} under a specific objective function, reached
deterministically from a census-derived seed via \cs{} \citep{dellaLibera2026b16}.

\paragraph{Two different goals.}
We insist on a clean distinction:

\begin{keybox}
\textbf{Goal A (Ensemble as Evaluator)}: Generate a distribution of
``reasonable'' plans from the geographic constraints alone.  Use this
distribution to evaluate any single plan---partisan, court-drawn, or
algorithmic---by its percentile position.

\textbf{Goal B (Algorithm as Canonical Plan)}: Compute a single plan that is
optimal by a stated criterion (minimum edge cut, maximum compactness, etc.)
and verifiably so.  The plan is canonical---reproducible from public data
by anyone.
\end{keybox}

These goals are complementary: the ensemble provides the reference distribution
against which the canonical plan's properties can be quantified.  But they
are not substitutes.  An ensemble plan is stochastic and unverifiable in the
sense that no particular sampled plan has a privileged claim.  The canonical
plan is deterministic and verifiable: given the 2020 census data and the
seat count, there is exactly one \ar{} map for each state.

\paragraph{The G-track.}
This paper establishes the methodology for the G-series of companion papers.
G.1 computes the percentile position of the \ar{} plan within published
ensemble results for six states.  G.2 focuses on partisan outcome distributions.
G.3 examines compactness distributions.  G.4 develops convergence diagnostics
($\rhat$, \ess{}, Hamming autocorrelation).  G.5 addresses mixing analysis
for the \recom{} chain itself.

\paragraph{Organisation.}
Section~\ref{sec:ensemble-methods} reviews ensemble methods.
Section~\ref{sec:percentile} introduces the percentile framework.
Section~\ref{sec:convergence} covers convergence diagnostics.
Section~\ref{sec:bridge} connects the diagnostics to \cs{}.
Section~\ref{sec:statutory} discusses statutory implications.
Section~\ref{sec:conclusion} concludes.
